\newcommand{\bs}{\symbol{'134}} % backslash in typewriter T1/T1 

\newcommand{\env}[1]{\texttt{#1}}

% For tex4ht packages splitindex and glossaries/glossaries-extra are not loaded, 
% makeindex is not invoked. 
% note how to write commands with optional arguments: 
% optional argument in {} like so: 
% \cmd[{[yy]}\{\dots\}\{\dots\}]{indexentry} 
\IfPackageLoadedTF{tex4ht}{
  \newcommand{\cmd}[2][]{\texttt{\bs{}#2#1}}
  \newcommand{\pkg}[1]{\texttt{#1}}% without indexing 
  \newcommand{\tool}[1]{\texttt{#1}}% without indexing 
  \newcommand{\param}[1]{\texttt{#1}}% without indexing 
}
{
  \newcommand{\cmd}[2][]{\texorpdfstring{\texttt{\bs{}#2#1}\sindex[cmd]{#2@\texttt{\bs{}#2}}}{\detokenize{\#2}}}
  \newcommand{\pkg}[1]{\texttt{#1}\sindex[pkg]{#1}} % TBD: this must be extracted 
  \newcommand{\tool}[1]{\texttt{#1}\sindex[tool]{#1}}% TBD: this must be extracted 
  \newcommand{\param}[1]{\texttt{#1}\sindex[param]{#1}}% TBD: this must be extracted 
}

% marking tools without making an entry in the index. 
% This is used for headlines/captions and in glossary 
\newcommand{\toolg}[1]{\texttt{#1}}

% TBD: the following 3 shall be replaced by \tool{...}
\newcommand{\pdflatex}{\tool{pdflatex}}
\newcommand{\lualatex}{\tool{lualatex}}
\newcommand{\xelatex}{\tool{xelatex}}


\newcommand{\texlive}{\TeX~Live}
\newcommand{\miktex}{MiKTeX}
\newcommand{\ltex}{L\TeX}% maybe also with \tool
\newcommand{\auctex}{AUC\TeX}

% TBD: clarify: \tool may be used in definition of acroyms, whereas \pkg may not. 
% Could be because \sindex does strange things 